\section{Introduction}
\label{sec:intro}

LLM serving systems batch many requests through one model. Each engine step generates
one token for every running request and may also process part of a new request's
prompt \citep{yu2022orca,kwon2023vllm,agrawal2024sarathiserve}. Service targets are
usually written in terms of time to first token (TTFT) and time per output token
\citep{zhong2024distserve}. A user feels both the wait for the first token and how
smoothly the rest stream. An operator who sees latency rise therefore needs to know
which phase grew and why. Past saturation the excess work is the same either way. The
scheduler decides whether it appears as waiting before the first token or as slower
steps, and its settings move it between the two.

This report asks under which conditions the dominant source of client-visible delay
changes, and which measurements can tell the causes apart. We
answer it for one setup, vLLM serving Qwen2.5-0.5B-Instruct on one RTX A4000. We use a
method designed to make negative results visible. Each experiment's predictions and
thresholds were published before its data were collected, runs were measured from
three independent views, and every prediction was scored afterward, including the
ones that failed.

Our contributions are:
\begin{itemize}
  \item \emph{A pre-registered, three-view measurement method} for LLM serving, with
    an opt-in per-step scheduler log that records how many tokens each engine step
    processed and when its results returned (Section~\ref{sec:method}).
  \item \emph{Where latency goes under load:} from C8 to C32, past saturation, the
    added latency appears as per-step time while the share before the first token
    falls, and a capacity argument bounds how far this extends
    (Section~\ref{sec:load}).
  \item \emph{The token budget as a median-versus-tail trade:} at C32 the budget sets
    the scale of the inter-token latency tail, and the smallest budget moves the delay
    into the queue (Section~\ref{sec:budget}).
  \item \emph{What sets the cost of a step:} context, measured directly for single
    chunks and consistent with a FLOP count, and what the step log could and could
    not show about scheduled tokens (Section~\ref{sec:stepcost}).
  \item \emph{A prediction ledger:} all 59 scored predictions and what the 19 that did
    not fully hold taught us (Section~\ref{sec:ledger}).
\end{itemize}
The code, configurations, protocols, and reports are public.\footnote{\url{https://github.com/pniyongabo/BottleneckShift}}
% src: ledger counts (40 confirmed, 4 mixed, 14 not confirmed, 1 not testable)
